\par\bigskip\pdfbookmark[0]{On the Connection. Fragment VI}{nexus-fragment-VI}\centerline{[VI.]}\medskip
\textsc{Theorem.} The number of classes, into which all properly primitive binary forms with negative determinant $-D$\begingroup\renewcommand{\thefootnote}{*)}\footnote{[are divided.]}\endgroup, is equal to
\[
\frac\pi4\times\surd D\times\text{Product of }\frac{p-1}p\frac{q+1}q\times\frac{rr-1}{rr}
\]
where
\begin{description}
\item[$p$] denotes all primes\begingroup\renewcommand{\thefootnote}{**)}\footnote{[odd.]}\endgroup\ for which $-D$ is a nonresidue,
\item[$q$] all primes\textsuperscript{**)} for which $-D$ is a residue,
\item[$r$] all primes\textsuperscript{**)} that divide $D$;
\end{description}
\[
=\frac{\frac\pi4\surd D\ \text{Product of }\frac{rr-1}{rr}}{1\pm\tfrac13\pm\tfrac15\text{ etc.}}
\]

in the denominator, the positive sign is put before fractions whose denominators have the form of nondivisors, and the negative sign before those whose denominators have the form of divisors of $xx+D$; fractions whose denominators are not coprime to $D$ are omitted altogether\begingroup\renewcommand{\thefootnote}{***)}\footnote{[If $m$ denotes all positive integers coprime to $2D$, and we use \textsc{Legendre}'s symbol as generalized by \textsc{Jacobi}, the above rule for determining the signs must be corrected as follows: in the preceding formula the denominator is
\[
1\pm\tfrac13\pm\tfrac15\text{ etc.}=\Sigma\pm\left(\frac{-D}m\right)\frac1m
\]
where the upper or lower sign is to be taken according as the number $m$ is a product of an even or odd number of primes (equal or unequal); in contrast, in the numerator of the following formula we have
\[
1\pm\tfrac13\pm\tfrac15\ldots=\Sigma\left(\frac{-D}m\right)\frac1m
\]
]}\endgroup.

\SourceLeaf{296}
\[
\begin{aligned}
&=\frac{2\surd D(1\pm\tfrac13\pm\tfrac15\ldots)}\pi\\
&=\frac{\operatorname{cotg}\theta\pm\operatorname{cotg}3\theta\pm\operatorname{cotg}5\theta\ldots\pm\operatorname{cotg}n\theta}{N:\surd D}
\end{aligned}
\]
putting $\theta=\frac\pi N$, $N=\left\{\begin{smallmatrix}1\\4\end{smallmatrix}\right\}D$, and taking for $n$ all numbers coprime to $D$, with the sign determined as above\begingroup\renewcommand{\thefootnote}{*)}\footnote{[See the note on this fragment further below.]}\endgroup.

For a positive determinant the number of classes will be\begingroup\renewcommand{\thefootnote}{**)}\footnote{[In the following formula $D$ denotes the positive determinant, and
\[
1\pm\tfrac13\pm\tfrac15\ldots=\Sigma\left(\frac Dm\right)\frac1m
\]
]}\endgroup
\[
=\frac{2\surd D(1\pm\tfrac13\pm\tfrac15\ldots)}{\log .T+U\surd D}
\]
Here $T,U$ denote the smallest values of the quantities $t,u$ that satisfy the equation $tt-Duu=1$
\[
=\frac{\log\sin\tfrac12\theta\pm\log\sin\tfrac32\theta\pm\log\sin\tfrac52\theta\text{ etc.}}{\log T+U\surd D}
\]
\BookRule{.30\textwidth}
